\chapter{Πειραματική αξιολόγηση}
\label{ch:evaluation}

% Εικόνα από το analysis/ αν υπάρχει, αλλιώς πλαίσιο-υπενθύμιση.
\newcommand{\genfig}[2][\linewidth]{%
  \IfFileExists{../figures/generated/#2}%
    {\includegraphics[width=#1]{#2}}%
    {\fbox{\parbox[c][4cm][c]{0.8\linewidth}{\centering\sffamily\small Θα παραχθεί από \code{thesis/analysis/}\\{\scriptsize\code{#2}}}}}}
\newcommand{\gentable}[1]{%
  \IfFileExists{../figures/generated/#1}%
    {\begin{adjustbox}{max width=\linewidth}\input{../figures/generated/#1}\end{adjustbox}}%
    {\fillin{πίνακας \code{#1} από το analysis/}}}

Το κεφάλαιο αξιολογεί το σύστημα ως προς τα τέσσερα ερευνητικά ερωτήματα της \cref{sec:rq}. Για κάθε πείραμα περιγράφονται ο στόχος, η διάταξη, η διαδικασία και τα μέτρα. Ακολουθούν τα αποτελέσματα και η συζήτησή τους. Όλοι οι αριθμοί, πίνακες και γραφήματα παράγονται από τα σενάρια ανάλυσης του φακέλου \code{thesis/analysis/} πάνω στα ακατέργαστα δεδομένα του \code{thesis/data/}, ώστε να μπορούν να αναπαραχθούν.

\review[!]{\textbf{Όλο το κεφάλαιο περιμένει μετρήσεις.} Το κείμενο της μεθοδολογίας είναι γραμμένο ως σχέδιο: όπου η πραγματική διαδικασία διαφέρει (αριθμός επαναλήψεων, αποστάσεις, χώρος), \textbf{διόρθωσε το κείμενο ώστε να περιγράφει ό,τι έκανες πραγματικά}, όχι ό,τι σχεδιάστηκε. Τα αποτελέσματα και η συζήτηση γράφονται από σένα, με βάση τα γραφήματα.}

\section{Πειραματική διάταξη}
\label{sec:exp-setup}

Όλα τα πειράματα έγιναν με Meta Quest~3 (λογισμικό συσκευής \fillin{έκδοση}) και την έκδοση της εφαρμογής που αντιστοιχεί στην ετικέτα \code{thesis-v2} του αποθετηρίου. Τα πειράματα πεδίου έγιναν στο \fillin{έργο / χώρος, περιγραφή, έκταση, αριθμός αισθητήρων}. Οι συντεταγμένες αναφοράς των σημείων προέρχονται από τοπογραφική αποτύπωση με \fillin{όργανο} με εκτιμώμενη ακρίβεια \fillin{x}\,mm. Η σύνδεση έγινε μέσω \fillin{Wi-Fi 5\,GHz / hotspot κινητού}.

\todo{Φωτογραφία ή κάτοψη του χώρου των πειραμάτων με τα σημεία ελέγχου και τα ανεξάρτητα σημεία.}

% ---------------------------------------------------------------------
\section{Ακρίβεια ευθυγράμμισης (ΕΕ1)}
\label{sec:exp-alignment}

\IfFileExists{../figures/generated/alignment_numbers.tex}{\input{../figures/generated/alignment_numbers.tex}}{}

\paragraph{Στόχος.} Να μετρηθεί πόσο απέχει ένα σημείο από τη σωστή του θέση μετά την ευθυγράμμιση, ως συνάρτηση του πλήθους και της γεωμετρίας των σημείων ελέγχου, και να ελεγχθεί πόσο καλά το προβλέπει η \cref{eq:predicted-error}.

\paragraph{Διάταξη.} Επιλέγονται 6--10 σημεία με γνωστές συντεταγμένες, κατανεμημένα στον χώρο σε αποστάσεις έως \fillin{x}\,m. Για κάθε επανάληψη, $k \in \{2, 3, 4\}$ από αυτά χρησιμοποιούνται για την ευθυγράμμιση και τα υπόλοιπα ως \emph{ανεξάρτητα σημεία ελέγχου}, που δεν συμμετέχουν στη λύση. Κάθε συνδυασμός επαναλαμβάνεται \fillin{n} φορές από την αρχή.

\paragraph{Εργαλείο μέτρησης.} Για το πείραμα προστέθηκε στην εφαρμογή η λειτουργία «Σημεία ελέγχου ακρίβειας» (\code{AlignmentAccuracy}). Ανοίγει μαζί με τη δοκιμή anchors από μία ροή, «Ξεκινάω μετρήσεις» (Μενού › Ευθυγράμμιση), που ζητά πρώτα τις συνθήκες και αν άλλαξε κάτι στο εργοτάξιο, και είναι κλειδωμένη στη λειτουργία Γραφείο. Ενεργοποιείται μόνο όταν υπάρχει ευθυγράμμιση, και τότε «παγώνει» την τρέχουσα λύση ως μια \emph{σειρά} μετρήσεων, με το πλήθος $k$ και τα αναγνωριστικά των σημείων ελέγχου, τη βάση $b$ και το RMS της λύσης. Όσο είναι ενεργή, ένα άγγιγμα σημείου του έργου \emph{δεν} τοποθετεί σημείο και δεν αλλάζει τη λύση. Η θέση του αγγίγματος μεταφέρεται στον κάναβο του έργου μέσω της αντίστροφης του μετασχηματισμού (\cref{eq:yaw,eq:translation}) και συγκρίνεται με τις συντεταγμένες του σημείου. Η διαφορά $(\Delta E, \Delta N, \Delta H)$ εμφανίζεται αμέσως στο πάνελ και γράφεται σε αρχείο CSV, μαζί με την απόσταση $D$ του σημείου από το κέντρο των σημείων ελέγχου και το προβλεπόμενο σφάλμα της \cref{eq:predicted-error}. Κάθε γραμμή χαρακτηρίζεται ως \emph{ανεξάρτητη} μέτρηση (σημείο εκτός της λύσης) ή ως μέτρηση \emph{σημείου ελέγχου} (σημείο της λύσης, ξανά-αγγιγμένο). Το δεύτερο μετρά την επαναληψιμότητα του αγγίγματος. Μια λανθασμένη μέτρηση αναιρείται επιτόπου με το κουμπί B, και μια νέα ευθυγράμμιση ξεκινά νέα σειρά. Τα αρχεία διατίθενται μέσω \code{GET /api/accuracy} και αναλύονται με το \code{analysis/alignment_accuracy.py}.

Η μέτρηση γίνεται στον κάναβο του έργου και όχι στον χώρο της συσκευής. Έτσι το σφάλμα αναλύεται σε συνιστώσες Ανατολής, Βορρά και ύψους, όπως σε μια τοπογραφική αποτύπωση, και είναι ανεξάρτητο από τη θέση της αρχής παρακολούθησης της συσκευής.

\paragraph{Διαδικασία.} Για κάθε σειρά: ευθυγράμμιση με $k$ σημεία· ενεργοποίηση της μέτρησης· άγγιγμα κάθε ανεξάρτητου σημείου μία φορά· άγγιγμα ενός σημείου ελέγχου \fillin{3} φορές για την εκτίμηση της επαναληψιμότητας· απενεργοποίηση. Οι συνθήκες φωτισμού δηλώνονται στο μενού και καταγράφονται σε κάθε γραμμή.

\paragraph{Μέτρα.} Το RMS της λύσης στα σημεία ελέγχου, όπως το δίνει η εφαρμογή. Το οριζόντιο σφάλμα $e_H = \sqrt{\Delta E^2 + \Delta N^2}$ και το κατακόρυφο $e_V = \Delta H$ στα ανεξάρτητα σημεία (διάμεσος, 95ο εκατοστημόριο). Η σχέση του $e_H$ με την απόσταση $D$. Ο λόγος μετρημένου προς προβλεπόμενο σφάλμα. Η επαναληψιμότητα του αγγίγματος: από τα επαναλαμβανόμενα αγγίγματα του ίδιου σημείου στην ίδια σειρά υπολογίζεται η συγκεντρωτική τυπική απόκλιση ανά άξονα,
\begin{equation}
  \hat\sigma_E^2 = \frac{\sum_j \sum_i (E_{ij} - \bar E_j)^2}{\sum_j (n_j - 1)},
\end{equation}
όπου $j$ τα σημεία και $i$ τα αγγίγματα, και ομοίως για $N$ και $H$. Η οριζόντια εκτίμηση $\hat\sigma = \sqrt{(\hat\sigma_E^2 + \hat\sigma_N^2)/2}$ συγκρίνεται με την τιμή $\sigma = \SI{2}{cm}$ που υποθέτει η \cref{eq:predicted-error}.

\begin{table}[htbp]
  \centering\small
  \caption{Ακρίβεια ευθυγράμμισης ανά πλήθος σημείων ελέγχου $k$ ($N$: ανεξάρτητες μετρήσεις).}
  \label{tab:alignment-by-k}
  \gentable{alignment_by_k.tex}
\end{table}

\begin{figure}[htbp]
  \centering
  \begin{subfigure}{0.56\linewidth}
    \genfig{alignment_error_vs_distance.pdf}
    \caption{Σφάλμα ως προς την απόσταση $D$. Διακεκομμένες: πρόβλεψη.}
    \label{fig:alignment-error}
  \end{subfigure}\hfill
  \begin{subfigure}{0.42\linewidth}
    \genfig{alignment_pred_vs_meas.pdf}
    \caption{Μετρημένο έναντι προβλεπόμενου.}
    \label{fig:alignment-pred}
  \end{subfigure}
  \caption{Οριζόντιο σφάλμα στα ανεξάρτητα σημεία, για $k = 2, 3, 4$.}
  \label{fig:alignment}
\end{figure}

\paragraph{Αποτελέσματα.}\mbox{}

\IfFileExists{../figures/generated/alignment_numbers.tex}{%
Από \AccChecks{} ανεξάρτητες μετρήσεις, το διάμεσο οριζόντιο σφάλμα ήταν \AccMedianCheckHcm\,cm και το 95ο εκατοστημόριο \AccPNinetyFiveCheckHcm\,cm. Η επαναληψιμότητα του αγγίγματος εκτιμήθηκε σε $\hat\sigma = \AccSigmaHcm$\,cm οριζόντια ανά άξονα και \AccSigmaVcm\,cm κατακόρυφα (\AccSigmaDof{} βαθμοί ελευθερίας).}{}
\todo{Σχόλιο στον \cref{tab:alignment-by-k} και στο \cref{fig:alignment}: αυξάνει το σφάλμα γραμμικά με $D$ και μειώνεται με τη βάση $b$; Υπερεκτιμά ή υποεκτιμά η πρόβλεψη (λόγος μετρ./προβλ.); Συμφωνεί το μετρημένο $\hat\sigma$ με τα \SI{2}{cm}; Τι σημαίνουν τα νούμερα για την πράξη (π.χ. πόσα σημεία, σε ποια απόσταση);}

% ---------------------------------------------------------------------
\section{Σταθερότητα και επανεντοπισμός anchors (ΕΕ2)}
\label{sec:exp-anchors}

\paragraph{Στόχος.} Να μετρηθεί αν τα spatial anchors ανά σημείο ελέγχου (\cref{sec:impl-anchors}) ξαναβρίσκονται σε επόμενες επισκέψεις, πόσο γρήγορα, και πόσο έχουν μετακινηθεί, υπό διαφορετικές συνθήκες.

\paragraph{Διαδικασία.} Τρία σημεία ελέγχου σε αποστάσεις \SIrange{5}{30}{m} μεταξύ τους. Την πρώτη ημέρα η ευθυγράμμιση γίνεται κανονικά με ενεργή τη «Δοκιμή πεδίου anchors», που καρφώνει ένα anchor σε κάθε σημείο. Ακολουθούν επισκέψεις μετά από 10 λεπτά, μία ημέρα και μία εβδομάδα. Σε κάθε επίσκεψη η εφαρμογή φορτώνει και επανεντοπίζει τα anchors (έως \SI{20}{s}) και καταγράφει το αποτέλεσμα. Στη συνέχεια ο χρήστης ακουμπά ξανά τα ίδια σημεία, και καταγράφεται η απόκλιση κάθε anchor από το νέο άγγιγμα. Οι επισκέψεις επαναλαμβάνονται υπό διαφορετικό φωτισμό (ηλιοφάνεια, συννεφιά, πρωί, μεσημέρι, βράδυ), μετά από αλλαγές στο εργοτάξιο, και μετά από απομάκρυνση \SI{50}{m} και επιστροφή χωρίς κλείσιμο της εφαρμογής. Οι συνθήκες δηλώνονται από τον χρήστη πριν από κάθε μέτρηση και καταγράφονται σε κάθε γραμμή.

\paragraph{Μέτρα.} Ποσοστό επιτυχούς επανεντοπισμού και διάμεσος χρόνος (γραμμές \code{load}/\code{manual}). Οριζόντια απόκλιση anchor--νέου αγγίγματος (διάμεσος, 95ο εκατοστημόριο) και απόλυτη κατακόρυφη απόκλιση (γραμμές \code{touch}). Τα μέτρα δίνονται ανά συνθήκη και ως συνάρτηση της ηλικίας του anchor. Η ανάλυση γίνεται με το \code{analysis/fieldtest_anchors.py}.

\begin{table}[htbp]
  \centering\small
  \caption{Επανεντοπισμός και απόκλιση anchors ανά συνθήκη.}
  \label{tab:anchors-by-condition}
  \gentable{anchors_by_condition.tex}
\end{table}

\begin{figure}[htbp]
  \centering
  \begin{subfigure}{0.49\linewidth}
    \genfig{anchors_offset_by_condition.pdf}
    \caption{Οριζόντια απόκλιση ανά συνθήκη.}
  \end{subfigure}\hfill
  \begin{subfigure}{0.49\linewidth}
    \genfig{anchors_offset_vs_age.pdf}
    \caption{Απόκλιση συναρτήσει ηλικίας anchor.}
  \end{subfigure}
  \caption{Σταθερότητα των spatial anchors.}
  \label{fig:anchors}
\end{figure}

\paragraph{Αποτελέσματα.}\mbox{}

\todo{Σχόλιο στον \cref{tab:anchors-by-condition} και στο \cref{fig:anchors}: σε ποιες συνθήκες τα anchors είναι αξιόπιστα; Αυξάνεται η απόκλιση με τον χρόνο; Τι σημαίνει αυτό για την αυτόματη επαναευθυγράμμιση;}

% ---------------------------------------------------------------------
\section{Απόδοση γραφικών (ΕΕ3)}
\label{sec:exp-performance}

\paragraph{Στόχος.} Να βρεθούν τα όρια φόρτου για σταθερό ρυθμό καρέ και η επίδραση του ρυθμού ανανέωσης και του foveated rendering.

\paragraph{Διάταξη.} Το εργαλείο OVR Metrics Tool καταγράφει τον ρυθμό καρέ, τους χρόνους CPU και GPU, τη θερμοκρασία και την μπαταρία σε αρχείο CSV. Τα σενάρια φόρτου είναι σωρευτικά:
(α) 50, 200 και 1000 σημεία·
(β) με ετικέτες, κορδέλες ιστορικού και ράβδους ορίων·
(γ) με μοντέλο 100.000, 300.000, 500.000 και 1.000.000 τριγώνων·
(δ) με χάρτη θερμότητας.
Κάθε σενάριο μετριέται για 2 λεπτά σε τρεις ρυθμίσεις: \SI{72}{Hz} χωρίς foveation, \SI{90}{Hz} χωρίς foveation, \SI{90}{Hz} με υψηλό foveation. Τα σημεία δημιουργούνται με το \code{fake_scenarios.py}, ώστε το σενάριο να είναι ίδιο σε κάθε επανάληψη.

\paragraph{Μέτρα.} Μέσος ρυθμός καρέ και 1\% χαμηλότερος, χρόνος καρέ GPU και CPU, ποσοστό καρέ που χάθηκαν, θερμοκρασία, κατανάλωση μπαταρίας ανά 10 λεπτά. Επιπλέον, ο χρόνος φόρτωσης μοντέλου (λήψη, ανάλυση, κατασκευή meshes) ως συνάρτηση του μεγέθους του αρχείου.

\begin{figure}[htbp]
  \centering
  \genfig[0.8\linewidth]{fps_vs_triangles.pdf}
  \caption{Ρυθμός καρέ συναρτήσει του πλήθους τριγώνων του μοντέλου, για τις τρεις ρυθμίσεις.}
  \label{fig:fps-triangles}
\end{figure}

\paragraph{Αποτελέσματα.}\mbox{}

\todo{Πού πέφτει ο ρυθμός κάτω από τα 72/90; Πόσο βοηθά το foveation; Επιβεβαιώνεται το όριο των 400.000 τριγώνων που χρησιμοποιεί η προετοιμασία μοντέλων;}

% ---------------------------------------------------------------------
\section{Προετοιμασία μοντέλων φωτογραμμετρίας}
\label{sec:exp-modelprep}

\paragraph{Στόχος.} Να ποσοτικοποιηθεί ο συμβιβασμός ποιότητας--απόδοσης της απλοποίησης (\cref{sec:impl-models}).

\paragraph{Διαδικασία.} Ένα μοντέλο φωτογραμμετρίας της κατασκευής απλοποιείται με το \code{model_prep.py} σε 1.000.000, 400.000, 200.000 και 100.000 τρίγωνα. Για κάθε εκδοχή μετρώνται: η απόσταση Hausdorff (μέση και μέγιστη) από το αρχικό πλέγμα με το pymeshlab \cite{cignoni1998metro}, το μέγεθος του πακέτου, ο χρόνος μεταφοράς και φόρτωσης, και ο ρυθμός καρέ (σενάριο (γ) της \cref{sec:exp-performance}).

\paragraph{Αποτελέσματα.}\mbox{}

\todo{Γράφημα: γεωμετρικό σφάλμα και ρυθμός καρέ συναρτήσει τριγώνων. Δικαιολογεί η καμπύλη την προεπιλογή των 400.000;}

% ---------------------------------------------------------------------
\section{Καθυστέρηση από άκρο σε άκρο}
\label{sec:exp-latency}

\paragraph{Στόχος.} Να μετρηθεί ο χρόνος από την αποστολή ενός σημείου ως την εμφάνισή του στη συσκευή.

\paragraph{Διαδικασία.} Ο εξυπηρετητής της συσκευής απαντά στο \code{POST /api/coordinates} μόλις βάλει τα σημεία στην ουρά του κύριου νήματος, πριν αυτά εμφανιστούν. Ο χρόνος απόκρισης του POST μετρά επομένως μόνο τη μεταφορά. Για τον χρόνο ως την εμφάνιση, ένα σενάριο Python στέλνει σημεία με μοναδικά αναγνωριστικά (1, 100 και 1000 ανά αίτημα) και αμέσως μετά ρωτά επανειλημμένα το \code{GET /api/coordinates}, που εξυπηρετείται από το κύριο νήμα, έως ότου εμφανιστούν όλα. Ο χρόνος από την αποστολή ως εκείνη την απάντηση είναι ένα άνω όριο της καθυστέρησης από άκρο σε άκρο, με ακρίβεια ενός καρέ (\SI{14}{ms} στα \SI{72}{Hz}). Η ίδια μέτρηση γίνεται και για τη διαδρομή push μέσω του \code{coordinate_server.py}. Δίκτυα: Wi-Fi \SI{5}{GHz} και hotspot κινητού.

\paragraph{Αποτελέσματα.}\mbox{}

\todo{Πίνακας: διάμεσος και 95ο εκατοστημόριο ανά μέγεθος αιτήματος και δίκτυο.}

% ---------------------------------------------------------------------
\section{Σύγκριση μεθόδων παρεμβολής}
\label{sec:exp-interpolation}

\paragraph{Στόχος.} Να συγκριθεί η γραμμική παρεμβολή σε τριγωνοποίηση Delaunay (η μέθοδος της εφαρμογής) με την IDW (η αρχική μέθοδος) στα πραγματικά δεδομένα του έργου.

\paragraph{Διαδικασία.} Διασταυρωμένη επικύρωση «αφήνοντας ένα έξω» (\en{leave-one-out}): για κάθε αισθητήρα, η μετακίνησή του εκτιμάται από τους υπόλοιπους με κάθε μέθοδο, και καταγράφεται η διαφορά από τη μετρημένη τιμή. Η ανάλυση γίνεται εκτός συσκευής, σε Python, με τις μετακινήσεις όλων των τοπογραφικών στόχων σε \fillin{n} ημερομηνίες. Μέθοδοι: γραμμική σε Delaunay, IDW με $p = 1, 2, 3$, και πλησιέστερος γείτονας ως βάση σύγκρισης.

\paragraph{Μέτρα.} RMSE και μέσο απόλυτο σφάλμα ανά μέθοδο, χωριστά για αισθητήρες στο εσωτερικό και στην περίμετρο της διάταξης.

\paragraph{Αποτελέσματα.}\mbox{}

\todo{Πίνακας RMSE/MAE. Επιβεβαιώνει η μέτρηση την επιλογή της γραμμικής παρεμβολής;}

% ---------------------------------------------------------------------
\section{Μελέτη χρηστών (ΕΕ4)}
\label{sec:exp-userstudy}

\paragraph{Στόχος.} Να διαπιστωθεί αν η επιτόπια απεικόνιση βελτιώνει την ταχύτητα και την ακρίβεια εργασιών επιθεώρησης σε σχέση με τη συμβατική ροή, και πώς αξιολογούν οι χρήστες τη χρηστικότητά της.

\paragraph{Συμμετέχοντες.} \fillin{N} συμμετέχοντες (\fillin{μηχανικοί / φοιτητές}), \fillin{x} με προηγούμενη εμπειρία σε VR/AR. Η συμμετοχή ήταν εθελοντική, με γραπτή συναίνεση (\cref{app:userstudy}).

\paragraph{Σχεδιασμός.} Σχεδιασμός εντός υποκειμένων (\en{within-subjects}) με δύο συνθήκες: \textbf{Α}, κάτοψη και λίστα αισθητήρων σε tablet· \textbf{Β}, η εφαρμογή. Η σειρά των συνθηκών αντισταθμίζεται: οι μισοί συμμετέχοντες ξεκινούν από την Α. Για να μην ωφελεί η εξοικείωση τη δεύτερη συνθήκη, χρησιμοποιούνται δύο ισοδύναμα σύνολα αισθητήρων. Πριν από τη συνθήκη Β κάθε συμμετέχων εξασκείται με το έργο εξάσκησης για \fillin{5} λεπτά.

\paragraph{Εργασίες.} (1) Εντοπισμός και άγγιγμα τριών συγκεκριμένων αισθητήρων. (2) Εύρεση των αισθητήρων σε κατάσταση συναγερμού και των θέσεών τους. (3) Εντοπισμός της περιοχής με τη μεγαλύτερη μετακίνηση τον τελευταίο μήνα.

\paragraph{Μέτρα.} Χρόνος ολοκλήρωσης και λάθη ανά εργασία. Χρηστικότητα με το SUS \cite{brooke1996sus} και φόρτος εργασίας με το NASA-TLX \cite{hart1988nasatlx}, και τα δύο μετά από κάθε συνθήκη. Συμπτώματα ναυτίας με το SSQ \cite{kennedy1993ssq} μετά τη συνθήκη Β. Ημιδομημένη συνέντευξη στο τέλος.

\paragraph{Ανάλυση.} Λόγω του μικρού δείγματος και των ζευγαρωτών μετρήσεων, οι συνθήκες συγκρίνονται με το μη παραμετρικό κριτήριο Wilcoxon \cite{wilcoxon1945}, με επίπεδο σημαντικότητας 0{,}05, και αναφέρεται μέγεθος επίδρασης.

\review[!]{Η μελέτη χρηστών χρειάζεται \textbf{έγκριση από την επιτροπή ηθικής} του ιδρύματος, αν το απαιτεί ο κανονισμός, και ρεαλιστικό αριθμό συμμετεχόντων. Αν γίνει πιλοτική μελέτη με 5--6 άτομα, δήλωσέ το ρητά και μην κάνεις στατιστικούς ελέγχους σημαντικότητας: δώσε περιγραφικά στατιστικά.}

\paragraph{Αποτελέσματα.}\mbox{}

\todo{Πίνακας χρόνων/λαθών ανά εργασία και συνθήκη, SUS και NASA-TLX, αποτελέσματα Wilcoxon, κύρια ευρήματα της συνέντευξης.}

% ---------------------------------------------------------------------
\section{Ποιότητα λογισμικού}
\label{sec:exp-quality}

Τα εργαλεία Python ελέγχονται με 127 unit tests (pytest). Καλύπτουν τους μετασχηματισμούς συντεταγμένων, την προετοιμασία μοντέλων και τη μεταφορά τους στον κάναβο των αισθητήρων, την αποστολή μοντέλων, τους πελάτες και το API του εξυπηρετητή push. Όλα περνούν στην έκδοση \code{thesis-v2}. Ο κώδικας C\# της εφαρμογής μεταγλωττίζεται χωρίς σφάλματα και χωρίς προειδοποιήσεις μεταγλωττιστή.

\todo{Προαιρετικά: κάλυψη κώδικα με \code{pytest --cov} και πίνακας ανά αρχείο.}
